% Generated from chapter-09.md - do not edit.

\chapter{Who Is Doing and Being?}%
\index{doing and being}\nobreak%

It can be helpful to have a simple model for understanding the different parts of ourselves and the different thoughts and needs that come from them. There are many ways of thinking about our inner selves, but the model that has been most useful to me is Parent, Adult, and Child from Transactional Analysis. It helps me sort out what is happening inside me and consider what thoughts might be most helpful.\index{Transactional Analysis}

\section{Meeting Your Inner Parent, Adult, and Child}

One of the ideas about our inner life comes from \textbf{Transactional Analysis}:

We have more than one ``voice'' inside us.

\emph{(Which explains a lot.)}

Transactional Analysis teaches that we carry within us three main parts:

\begin{itemize}
\item \textbf{The Parent}
\item \textbf{The Adult}
\item \textbf{The Child}
\end{itemize}

Each of these voices influences how we think, feel, and respond to life.

When I learned this, a lot of things suddenly made sense.

\emph{(Including some of my more confusing moments.)}

\section{The Parent}%
\index{parent}\nobreak%

The Parent has two main sides.

One is the \textbf{Critical Parent}---the part that corrects, warns, scolds, gives advice, and says a lot of ``shoulds.''

\emph{(It can be very persuasive.)}

The other is the \textbf{Nurturing Parent}---the part that says:

\begin{itemize}
\item ``I love you.''
\item ``I care about you.''
\item ``Come here, let me hold you.''
\end{itemize}

Both of these developed for good reasons.

They were meant to protect us and guide us.

But when the Critical Parent inside us runs the show . . . 

\textbf{Better Thoughts tend to disappear rather quickly.}

\section{The Child}%
\index{child}\nobreak%

The Child is the emotional part of us.

The Child can be:

\begin{itemize}
\item playful and joyful
\item sensitive and tender
\item fearful and anxious
\item people-pleasing
\item rebellious
\item free and spontaneous
\end{itemize}

The Child brings life, fun, joy, color, creativity, and feeling.

But the Child can also become overwhelmed, frightened, and uncertain.

When the Child is in charge and scared, our thinking often becomes small, tight, and worried.

\emph{(Everything suddenly feels like a very big deal.)}

\section{The Adult}%
\index{adult}\nobreak%

The Adult is the steady, present, thoughtful part of us.

The Adult:

\begin{itemize}
\item thinks clearly, calmly
\item makes choices
\item responds to what is actually happening now
\item looks for facts and possibilities
\item solves problems
\item regulates emotion
\item brings calm leadership inside
\end{itemize}

This is the part of me I most want to strengthen.

Because when my Adult is present . . . 

\textbf{Better Thoughts appear much more naturally.}

\section{Inviting the Adult Forward}

One of my ongoing goals is to increase my Adult's capacity to be in charge---especially when my Child is anxious.

I do that by saying \textbf{Adult things on purpose}.

Things like:

\begin{itemize}
\item ``I can handle whatever comes along.''
\item ``I can stop and think about what ought to be done.''
\item ``I'm a good problem solver.''
\item ``I can meet obstacles without fear.''
\item ``This moment is manageable.''
\item ``I choose my next step.''
\end{itemize}

These are not just positive statements.

They are \textbf{Adult statements}---steady, capable, grounded in the present.

They remind my nervous system who is in charge.

\emph{(A very helpful reminder.)}

\section{Using the Body to Support the Adult}

My Adult also uses the body to stay grounded.

Sometimes my Adult says:

\begin{itemize}
\item ``Stand up straight.''
\item ``Take a slow breath.''
\item ``Feel your feet.''
\item ``Look around the room.''
\item ``Notice where you are sitting.''
\end{itemize}

These are simple cues that bring me back into the here and now---into the Adult.

From that place, my thinking clears.

\section{The Adult's Three Helpful Questions}

When I feel anxious or uncertain, my Adult asks:

\begin{itemize}
\item What are the facts?
\item What is really happening right now?
\item What needs to be done?
\item What are the possibilities for getting there?
\end{itemize}

This moves me out of fear and into wise action.

\emph{(A much more comfortable place to be.)}

\section{Helping the Child with the Adult}

When my Child is scared, my Adult gently steps in and says:

\begin{itemize}
\item ``I see you.''
\item ``I know you're afraid.''
\item ``You don't have to do this alone.''
\item ``I am here now.''
\item ``I will handle the thinking and planning.''
\item ``You can rest with me.''
\end{itemize}

This is where anxiety begins to soften.

The Child doesn't disappear.\\
The Parent doesn't disappear.

But the Adult becomes the leader.

And from that place, Better Thoughts grow naturally.

\section{Meeting the Needs of My Inner Family}

\emph{(Parent, Adult, and Child)}

Sometimes I ask myself, \emph{``What's a better thought?''}

And sometimes the better question is:

\textbf{``Which part of me has a need right now?''}

I carry within me a Parent, an Adult, and a Child.

Not three different people---just three ways of responding to life.

\emph{(A small internal committee.)}

If I ignore their needs, I feel restless, critical, impulsive, or off course.

If I listen, something settles.

\section{My Child Needs}

My Child needs comfort, affection, play, and reassurance.

She needs to feel safe.

She needs someone to say, \emph{``You're okay. You're loved.''}

If I'm reaching for something I don't really want---extra food, mindless television, or reassurance from someone who isn't available---I sometimes ask:

Is my Child lonely? Tired? Unseen?

Often she just needs:

\begin{itemize}
\item a nap
\item a kind word, some love and attention
\item or a memory of being loved
\end{itemize}

\emph{(Sometimes a nap solves more than expected.)}

\section{My Parent Needs}

My Parent carries my values.

She wants me to be good, responsible, loving, and spiritually aligned.

When I feel guilty or disappointed in myself, I ask:

Is my Parent afraid I'm going off track?

Sometimes she needs reassurance that I am still a good person---even when I'm imperfect.

Sometimes she needs to soften and remember that kindness works better than criticism.

\emph{(It usually does.)}

\section{My Adult Needs}

My Adult needs information, clarity, and calm.

When I'm anxious about money, health, or decisions, my Adult simply needs:

\begin{itemize}
\item facts
\item time to think
\item a deep breath
\end{itemize}

When my Adult is present, I remember:

I can choose.\\
I can pause.\\
I can respond instead of react.

\section{Pause for a Moment}

Right now . . . 

Which part of you is speaking?

The Child?\\
The Parent?\\
The Adult?

And what might that part need?

No fixing required.

Just a little understanding.

\section{A Kinder Question}

Instead of asking, \emph{``What's wrong with me?''}

I can ask:

\textbf{``Which part of me is speaking---and what does it need?''}

Meeting my own needs is not selfish.

It is wise stewardship of my inner life.

And here is the beautiful surprise:

When my inner Parent feels calm,\\
my inner Child feels safe,\\
and my Adult feels clear . . . 

I am much more likely to feel happy.

Not giddy.\\
Not perfect.

But peacefully, steadily well.

\section{Examples of Meeting My Own Needs}

When I lie down for a nap instead of pushing through \ensuremath{\rightarrow} I am caring for my Child.

When I forgive myself for a mistake \ensuremath{\rightarrow} I am softening my Parent.

When I check my bank balance and review the facts \ensuremath{\rightarrow} I am strengthening my Adult.

When I call a friend and ask for listening \ensuremath{\rightarrow} I am honoring my Child's need to be heard.

When I pray and then pause before acting \ensuremath{\rightarrow} I am giving my Adult space.

\section{A Better Thought}

A very good Better Thought for me is simply this:

\textbf{``My Adult is here---and that is enough for this moment.''}

I have a feeling that this idea holds a great deal.

I do have a friend who believes her thoughts should go wherever they want to go---and is willing to suffer through some very discouraging ones.

But I've discovered something about myself:

\textbf{I want my thoughts to lead me toward something kinder . . .  and happier.}

\emph{(That seems like a reasonable preference.)}

\emph{------------------------------------------}
